\chapter{เกณฑ์การคำนวณดัชนีคุณภาพอากาศของประเทศไทย}
\label{app:appendixB}

\section{ตารางเปรียบเทียบระดับความเข้มข้น PM2.5 กับค่าดัชนีคุณภาพอากาศ (AQI)}
อ้างอิงตามประกาศกรมควบคุมมลพิษ (พ.ศ. 2566):

\begin{table}[H]
\centering\small
\begin{tabularx}{\textwidth}{c c p{0.25\textwidth} X}
\toprule
\textbf{ค่า AQI} & \textbf{PM2.5 ($\mu g/m^3$)} & \textbf{ระดับคุณภาพอากาศ} & \textbf{ความหมายและข้อแนะนำด้านสุขภาพ} \\
\midrule
0 -- 25 & 0 -- 15.0 & ฟ้า (ดีมาก) & คุณภาพอากาศดีมาก เหมาะสำหรับกิจกรรมกลางแจ้ง \\
26 -- 50 & 15.1 -- 25.0 & เขียว (ดี) & คุณภาพอากาศดี สามารถทำกิจกรรมกลางแจ้งได้ตามปกติ \\
51 -- 100 & 25.1 -- 37.5 & เหลือง (ปานกลาง) & ประชาชนทั่วไปทำกิจกรรมได้ ผู้มีปัญหาสุขภาพควรลดกิจกรรมหนัก \\
101 -- 200 & 37.6 -- 75.0 & ส้ม (เริ่มมีผลกระทบ) & ควรลดกิจกรรมกลางแจ้ง สวมหน้ากากป้องกันฝุ่น N95 \\
$>$ 200 & $>$ 75.0 & แดง (มีผลกระทบต่อสุขภาพ) & ทุกคนควรหลีกเลี่ยงกิจกรรมกลางแจ้ง งดการออกกำลังกาย \\
\bottomrule
\end{tabularx}
\end{table}
